\documentclass[12pt,a4paper]{article}
\usepackage{amsmath,amssymb,amsthm}
\usepackage{hyperref}
\usepackage{geometry}
\geometry{margin=2.5cm}

\newtheorem{theorem}{Theorem}[section]
\newtheorem{lemma}[theorem]{Lemma}
\newtheorem{corollary}[theorem]{Corollary}
\newtheorem{conjecture}[theorem]{Conjecture}
\theoremstyle{definition}
\newtheorem{definition}[theorem]{Definition}
\theoremstyle{remark}
\newtheorem{remark}[theorem]{Remark}

\title{The Euler--Mascheroni Constant in Recognition Science:\\
$\gamma$ as a Ledger-Harmonic Quantity}

\author{Jonathan Washburn\\
Recognition Science Research Institute\\
Austin, Texas\\
\texttt{washburn.jonathan@gmail.com}}

\date{March 2026}

\begin{document}
\maketitle

\begin{abstract}
We establish the role of the Euler--Mascheroni constant
$\gamma = \lim_{n\to\infty}(H_n - \ln n) \approx 0.5772$ in
Recognition Science~\cite{RS_Axioms}.  We prove the bounds
$1/2 < \gamma < 2/3$ from elementary estimates on the harmonic
series.  The constant appears in RS as the harmonic-series accumulation
of ledger costs: $H_n = \sum_{k=1}^n 1/k$ is the cumulative cost of
$n$ successive recognition events at integer scales, and $\gamma$
measures the asymptotic excess over the continuum approximation
$\ln n$.  A closed-form $\varphi$-expression for $\gamma$ remains
open, as does the irrationality of $\gamma$.
\end{abstract}

%% ============================================================
\section{Recognition Science Framework}
\label{sec:framework}
%% ============================================================

The Euler--Mascheroni constant $\gamma$ arises in RS as the
asymptotic discrete excess of harmonic ledger costs.
Its role is grounded in the following structure.

\begin{definition}[RS Cost Functional]
For $x > 0$: $J(x) = \tfrac{1}{2}(x + x^{-1}) - 1$.
\end{definition}

\noindent\textbf{Axiom A1 (Normalization):} $J(1) = 0$.\\
\textbf{Axiom A2 (Recognition Composition Law, RCL):}
$J(xy)+J(x/y)=2J(x)J(y)+2J(x)+2J(y)$ for all $x,y>0$.\\
\textbf{Axiom A3 (Calibration):} $J''_{\log}(0) = 1$, where
$J_{\log}(t) := J(e^t)$.

\begin{theorem}[Cost Uniqueness]
\label{thm:unique}
The continuous solution to axioms \textup{(A1)--(A3)} is
unique: $J(x) = \tfrac{1}{2}(x + x^{-1}) - 1$.
\end{theorem}

\begin{proof}[Proof sketch]
Set $H(t) = J_{\log}(t)+1$.  Axiom A2 transforms into the d'Alembert
equation $H(s+t)+H(s-t)=2H(s)H(t)$.  By the Acz\'el classification
theorem~\cite{Aczel1966}, the unique continuous solution with $H(0)=1$
and $H''(0)=1$ is $H(t)=\cosh t$, giving
$J(x)=\tfrac{1}{2}(x+x^{-1})-1$.
\end{proof}

\noindent In RS, successive recognition events at integer scales $k=1,
2,\ldots,n$ carry cumulative cost $H_n = \sum_{k=1}^n 1/k$.
The continuum approximation of this sum is $\ln n$.  The
Euler--Mascheroni constant $\gamma = \lim_{n\to\infty}(H_n - \ln n)$
is therefore the irreducible cost of discreteness: the amount by
which the integer-step RS ledger permanently exceeds its smooth
continuum counterpart.

%% ============================================================
\section{Introduction}
\label{sec:intro}
%% ============================================================

The Euler--Mascheroni constant
$\gamma = \lim_{n\to\infty}(H_n - \ln n) \approx 0.5772$ appears
throughout mathematics and physics.  In number theory, it arises in
Mertens' theorems on prime distribution and in the analysis of
$\zeta(s)$ near $s = 1$~\cite{Lagarias}.  In quantum field theory,
$\gamma$ emerges in dimensional regularization and the
$\overline{\mathrm{MS}}$ scheme~\cite{Collins}.  Its irrationality
remains unproven~\cite{Papanikolaou}.

\section{Definition and Bounds}
\label{sec:bounds}

\begin{definition}[Euler--Mascheroni Constant]
\label{def:gamma}
\begin{equation}
  \gamma = \lim_{n\to\infty}\left(\sum_{k=1}^n \frac{1}{k}
  - \ln n\right) = \lim_{n\to\infty}(H_n - \ln n).
\end{equation}
\end{definition}

\begin{lemma}[Monotonicity]
\label{lem:mono}
The sequence $a_n = H_n - \ln n$ is strictly decreasing and
converges to $\gamma$.  The sequence $b_n = H_n - \ln(n+1)$ is
strictly increasing and converges to $\gamma$ from below.
\end{lemma}

\begin{proof}
For $a_n$: $a_{n+1} - a_n = 1/(n+1) - \ln((n+1)/n) < 0$ since
$\ln(1 + 1/n) > 1/(n+1)$ for $n \geq 1$.  For $b_n$: similarly
$b_{n+1} - b_n > 0$.  Both sequences are bounded and monotone,
hence converge to $\gamma$.
\end{proof}

\begin{theorem}[Certified Bounds]
\label{thm:bounds}
\begin{enumerate}
  \item $\gamma > 0$.
  \item $\gamma < 2/3$.
  \item $1/2 < \gamma < 2/3$.
  \item $\gamma \neq 0$.
\end{enumerate}
\end{theorem}

\begin{proof}
(1)~From Lemma~\ref{lem:mono}, $b_1 = 1 - \ln 2 \approx 1 - 0.693
= 0.307 > 0$, and $\gamma > b_1 > 0$.

(2)~Since $a_n = H_n - \ln n$ decreases to $\gamma$, we have
$\gamma < a_n$ for all $n$.  Take $n = 10$:
$H_{10} = 1 + \tfrac{1}{2} + \cdots + \tfrac{1}{10}
= \tfrac{7381}{2520} \approx 2.9290$ and $\ln 10 \approx 2.3026$,
giving $a_{10} \approx 0.6264 < 2/3$.  Hence $\gamma < 2/3$.

(3)~Take $n = 8$: $b_8 = H_8 - \ln 9$.
$H_8 = \tfrac{761}{280} \approx 2.7179$ and $\ln 9 \approx 2.1972$,
so $b_8 \approx 0.5207 > 1/2$, giving $\gamma > 1/2$.

(4)~Immediate from (1): $\gamma > 0 \neq 0$.
\end{proof}

\begin{remark}
All four bounds follow from elementary computation; no asymptotic
expansions are required.
\end{remark}

\begin{remark}
The bounds $1/2 < \gamma < 2/3$ are tight enough for all RS
applications: $\gamma = 0.5772\ldots$ lies well within the interval.
\end{remark}

\section{Role in Recognition Science}
\label{sec:role}

In RS, the harmonic series $H_n = \sum_{k=1}^n 1/k$ represents the
cumulative cost of $n$ successive recognition events at consecutive
integer scales.  The logarithmic approximation $\ln n$ captures the
smooth (continuum) part; $\gamma$ measures the purely discrete
remainder---the irreducible cost of discreteness itself.

\begin{remark}[Ledger interpretation]
\label{rem:ledger}
In RS, the discrete cost of $n$ successive recognition events at
integer scales is $H_n = \sum_{k=1}^n 1/k$, while the continuum
approximation is $\ln n$.  The asymptotic excess
\begin{equation}
  \gamma = \lim_{n\to\infty}(H_n - \ln n)
\end{equation}
is therefore the \emph{irreducible cost of discreteness}: the amount
by which the integer-step ledger permanently exceeds its smooth
approximation.
\end{remark}

\section{Physical Occurrences in RS}
\label{sec:physical}

\begin{remark}[Running Coupling]
\label{thm:running}
In the RS treatment of coupling evolution, the discrete sum
$\sum_{k=1}^n 1/k$ over recognition scales differs from the
continuum integral $\int_1^n dx/x = \ln n$ by exactly $\gamma$ in
the limit.  This mirrors QFT: $\gamma$ appears in dimensional
regularization and $\overline{\mathrm{MS}}$ beta functions.  The
occurrence of $\gamma$ in RS coupling evolution is a structural
observation; a derivation from the RS Hamiltonian is ongoing work.
\end{remark}

\begin{remark}[Harmonic Structure of the $\varphi$-Ladder]
\label{thm:ladder-harmonic}
The $\varphi$-ladder $m_n = m_0 \varphi^n$ has a natural harmonic
weighting.  The integer-scale sum $H_n$ diverges as $\ln n + \gamma$,
while $\sum_{k=1}^n \varphi^{-k} \to \varphi$.  Thus $\gamma$
measures the correction when integer-scale counting replaces
$\varphi$-scale counting.  Both $\gamma \approx 0.5772$ and
$\ln\varphi \approx 0.4812$ are positive universal constants in RS.
\end{remark}

\begin{remark}[Renormalization Excess]
\label{thm:renorm}
In RS renormalization, sums $\sum_{k=1}^\Lambda 1/k$ are regularized
by subtracting $\ln\Lambda$.  The finite remainder is $\gamma$,
independent of the cutoff.  Thus $\gamma$ is the universal discrete
excess in RS renormalization.  This is a direct consequence of the
definition of $\gamma$ (Definition~\ref{def:gamma}), not an
independent theorem.
\end{remark}

\section{Connection to the $\varphi$-Ladder}
\label{sec:phi-ladder}

The $\varphi$-ladder $m_n = m_0 \varphi^n$ yields $H_n$ when counting
rungs with cost $1/k$ per rung.

\begin{remark}[Discrete--Continuum Gap]
\label{thm:gap}
$\gamma$ is the asymptotic gap between the discrete harmonic sum
$H_n$ and the continuum approximation $\ln n$.  It quantifies the
irreducible discreteness of the $\varphi$-ladder when enumerated by
integer rungs.  This follows directly from
Definition~\ref{def:gamma} and the monotonicity
Lemma~\ref{lem:mono}.
\end{remark}

\begin{corollary}
The $\varphi$-ladder, when counted harmonically at integer scales,
accumulates an excess $\gamma$ over the logarithmic continuum.  This
excess is universal and cutoff-independent.
\end{corollary}

\section{Comparison: $\gamma$ and $\ln\varphi$}
\label{sec:comparison}

Two fundamental constants in RS are $\gamma \approx 0.5772$ and
$\ln\varphi \approx 0.4812$.  Both are positive and less than $1$;
the ratio $\gamma/\ln\varphi \approx 1.20$ is close to $6/5$.

\begin{remark}[Distinct Roles]
$\ln\varphi$ is the minimum $J$-cost (bit cost) of a non-trivial
recognition event; $\gamma$ is the asymptotic discrete excess of the
harmonic series.  The former is $\varphi$-derived; the latter arises
from integer-scale harmonic counting.  Both are positive constants
less than $1$; their ratio $\gamma/\ln\varphi \approx 1.20$ is close
to $6/5$ (open conjecture).
\end{remark}

\begin{remark}
Whether $\gamma = (6/5)\ln\varphi$ or a similar closed form holds
remains open.
\end{remark}

\section{Predictions}
\label{sec:predictions}

\begin{conjecture}[Closed Form]
$\gamma = f(\varphi, \zeta(2), \zeta(3), \ldots)$ for some function $f$.
\end{conjecture}

\begin{conjecture}[Irrationality]
$\gamma$ is irrational.
\end{conjecture}

\begin{conjecture}[RS Prediction]
In RS renormalization with harmonic sums, the finite remainder is
$\gamma$ (or a rational multiple).  Testable in coupling derivations.
\end{conjecture}

\begin{remark}
All conjectures are open; bounds $1/2 < \gamma < 2/3$ are certified.
\end{remark}

\section{Conclusion}
\label{sec:conclusion}

We have established the role of the Euler--Mascheroni constant
$\gamma \approx 0.5772$ in Recognition Science.  The constant is
certified within $1/2 < \gamma < 2/3$ with full proofs.  In RS,
$\gamma$ arises as the asymptotic discrete-cost excess of the
harmonic series.  It appears in three contexts: (i)~running of
coupling constants; (ii)~harmonic structure of the
$\varphi$-ladder; (iii)~renormalization as the universal finite
remainder.  We compared $\gamma$ with $\ln\varphi \approx 0.481$:
both are fundamental constants less than $1$ with distinct roles.  A
$\varphi$-expression for $\gamma$ and its irrationality remain open.

\begin{thebibliography}{9}

\bibitem{RS_Axioms}
J.~Washburn, ``The Algebra of Reality: A Recognition Science Derivation
of Physical Law,'' \emph{Axioms} \textbf{15}(2), 90 (2026).
\texttt{doi:10.3390/axioms15020090}

\bibitem{Aczel1966}
J.~Acz\'el, \emph{Lectures on Functional Equations and Their
Applications} (Academic Press, 1966).

\bibitem{Lagarias}
J.~C.~Lagarias, ``Euler's Constant: Euler's Work and Modern
Developments,'' Bull.\ AMS \textbf{50}, 527 (2013).

\bibitem{Collins}
J.~C.~Collins, \textit{Renormalization}, Cambridge University Press,
Cambridge (1984).

\bibitem{Papanikolaou}
T.~Papanikolaou, ``Euler's constant: irrationality and
transcendence,'' in \textit{Number Theory and Applications},
Kluwer, Dordrecht (1989).
\end{thebibliography}

\end{document}
